% \section*{NeurIPS Paper Checklist}

% \begin{enumerate}

% \item {\bf Claims}
%     \item[] Question: Do the main claims made in the abstract and introduction accurately reflect the paper's contributions and scope?
%     \item[] Answer: \answerYes{}
%     \item[] Justification: The abstract and Section~\ref{sec:intro} state four contributions (benchmark, red-line rubric taxonomy, per-stage trajectory finding, text-only ablation), each of which corresponds one-to-one to a main-paper section (\S\ref{sec:design}, \S\ref{sec:design-evaluation}, \S\ref{sec:exp-trajectory}, \S\ref{sec:exp-ablation}). The claimed top score of $0.758$ and the per-stage reversal are reported in Table~\ref{tab:leaderboard} and Figure~\ref{fig:trajectory}.

% \item {\bf Limitations}
%     \item[] Question: Does the paper discuss the limitations of the work performed by the authors?
%     \item[] Answer: \answerYes{}
%     \item[] Justification: Section~\ref{sec:limits} has a dedicated paragraph each on the mock-backend trade-off, the breadth-depth choice at 100 tasks, single-run main results with an Appendix~D variance bound, and checker-coverage honesty.

% \item {\bf Theory assumptions and proofs}
%     \item[] Question: For each theoretical result, does the paper provide the full set of assumptions and a complete (and correct) proof?
%     \item[] Answer: \answerNA{}
%     \item[] Justification: The paper contributes a benchmark and empirical findings; there are no theorems or formal proofs.

% \item {\bf Experimental result reproducibility}
%     \item[] Question: Does the paper fully disclose all the information needed to reproduce the main experimental results of the paper to the extent that it affects the main claims and/or conclusions of the paper (regardless of whether the code and data are provided or not)?
%     \item[] Answer: \answerYes{}
%     \item[] Justification: Every task ships with a \texttt{task.py} defining deterministic stage hooks and checkers (Section~\ref{sec:design-evaluation}); the harness is Dockerised and uses deterministic seeding (Section~\ref{sec:protocol}); we release all 100 tasks, the full harness, and 700 execution traces. Model-side sampling follows provider defaults, which we document in Appendix~H.

% \item {\bf Open access to data and code}
%     \item[] Question: Does the paper provide open access to the data and code, with sufficient instructions to faithfully reproduce the main experimental results, as described in supplemental material?
%     \item[] Answer: \answerYes{}
%     \item[] Justification: The benchmark tasks, mock-backend harness, checker library, and 700 per-task execution traces are released under a CC-BY-NC-4.0 data license and MIT code license; supplemental material includes the Docker image SHA, environment specification, and per-task run scripts (Appendix~H).

% \item {\bf Experimental setting/details}
%     \item[] Question: Does the paper specify all the training and test details (e.g., data splits, hyperparameters, how they were chosen, type of optimizer) necessary to understand the results?
%     \item[] Answer: \answerYes{}
%     \item[] Justification: Section~\ref{sec:protocol} specifies the seven evaluated models, the shared agent framework and tool schema, the 80-turn per-stage budget, the single-run protocol with the Appendix~D $k=3$ subset for variance, deterministic backend seeding, and the token-accounting convention (cache merged into input). No model fine-tuning is performed.

% \item {\bf Experiment statistical significance}
%     \item[] Question: Does the paper report error bars suitably and correctly defined or other appropriate information about the statistical significance of the experiments?
%     \item[] Answer: \answerYes{}
%     \item[] Justification: The main results (Table~\ref{tab:leaderboard}) are single-run scores; we bound run-to-run variance on a stratified 20-task $k=3$ subset in Appendix~D. The per-stage trajectory analysis of \S\ref{sec:exp-trajectory} uses a paired bootstrap over 73 three-stage tasks (10,000 resamples) with 95\% confidence intervals, taking task sampling as the randomness source.

% \item {\bf Experiments compute resources}
%     \item[] Question: For each experiment, does the paper provide sufficient information on the computer resources (type of compute workers, memory, time of execution) needed to reproduce the experiments?
%     \item[] Answer: \answerYes{}
%     \item[] Justification: All model inference uses hosted APIs; the client side runs inside a per-task Docker container (2 CPU, 4\,GB RAM) on commodity cloud instances. Table~\ref{tab:leaderboard} reports wall-clock time per full single-run sweep per model; Appendix~J reports total input and output tokens and estimated API spend per full sweep.

% \item {\bf Code of ethics}
%     \item[] Question: Does the research conducted in the paper conform, in every respect, with the NeurIPS Code of Ethics \url{https://neurips.cc/public/EthicsGuidelines}?
%     \item[] Answer: \answerYes{}
%     \item[] Justification: All task content was authored by the co-authors or drawn from publicly available professional scenarios with personally identifying information replaced by fictional personas; no real-world user data, medical records, or proprietary corporate data were used. The mock-backend harness runs in isolation and cannot affect real-world services.

% \item {\bf Broader impacts}
%     \item[] Question: Does the paper discuss both potential positive societal impacts and negative societal impacts of the work performed?
%     \item[] Answer: \answerYes{}
%     \item[] Justification: Section~\ref{sec:limits} discusses both: positive impact through falsifiable rule-based evaluation of coworker agents before professional deployment, and negative impact through lowering the friction to deploying agents in professional workflows. We note that no production action channel is released and the harness is development-only.

% \item {\bf Safeguards}
%     \item[] Question: Does the paper describe safeguards that have been put in place for responsible release of data or models that have a high risk for misuse (e.g., pre-trained language models, image generators, or scraped datasets)?
%     \item[] Answer: \answerYes{}
%     \item[] Justification: \bench{} releases no trained models. The released artifacts are tasks, checkers, mock-backend harness, and execution traces; the harness cannot invoke real email, calendar, knowledge-base, or spreadsheet services and therefore carries no direct misuse vector. Execution traces are redacted of any API keys used during evaluation.

% \item {\bf Licenses for existing assets}
%     \item[] Question: Are the creators or original owners of assets (e.g., code, data, models), used in the paper, properly credited and are the license and terms of use explicitly mentioned and properly respected?
%     \item[] Answer: \answerYes{}
%     \item[] Justification: The harness uses GreenMail (Apache-2.0) and Radicale (GPL-3.0); the knowledge-base and spreadsheet mock services are built against the providers' official public APIs. Credential-flow scripts are adapted from prior work credited in the supplementary material. All multimodal artifacts used in tasks were authored by the co-authors or drawn from CC-licensed sources; provenance is listed in Appendix~B.

% \item {\bf New assets}
%     \item[] Question: Are new assets introduced in the paper well documented and is the documentation provided alongside the assets?
%     \item[] Answer: \answerYes{}
%     \item[] Justification: The benchmark, harness, and checker library are released with a Datasheet (Appendix~A), per-task documentation (Appendix~B), and a reproducibility checklist (Appendix~H). Anonymised URLs for the task repository and execution traces are provided in the supplemental material.

% \item {\bf Crowdsourcing and research with human subjects}
%     \item[] Question: For crowdsourcing experiments and research with human subjects, does the paper include the full text of instructions given to participants and screenshots, if applicable, as well as details about compensation (if any)?
%     \item[] Answer: \answerNA{}
%     \item[] Justification: \bench{} is authored by the co-authors and does not use crowdsourced annotation or human subjects. No human-generated trajectory labels are collected.

% \item {\bf Institutional review board (IRB) approvals or equivalent for research with human subjects}
%     \item[] Question: Does the paper describe potential risks incurred by study participants, whether such risks were disclosed to the subjects, and whether Institutional Review Board (IRB) approvals (or an equivalent approval/review based on the requirements of your country or institution) were obtained?
%     \item[] Answer: \answerNA{}
%     \item[] Justification: No human subjects are involved in the research.

% \item {\bf Declaration of LLM usage}
%     \item[] Question: Does the paper describe the usage of LLMs if it is an important, original, or non-standard component of the core methods in this research? Note that if the LLM is used only for writing, editing, or formatting purposes and does \emph{not} impact the core methodology, scientific rigor, or originality of the research, declaration is not required.
%     \item[] Answer: \answerYes{}
%     \item[] Justification: LLMs are the object of evaluation, not a component of benchmark construction. The seven evaluated models are named in Section~\ref{sec:protocol}; benchmark scoring is entirely rule-based (no LLM-as-judge, Section~\ref{sec:design-evaluation}). LLM assistance was used during manuscript writing and editing only, in line with NeurIPS policy.

% \end{enumerate}
